\documentclass[a4paper,10pt]{article}
\usepackage[margin=22mm]{geometry}
\usepackage{amsmath}
\usepackage{lipsum}
\usepackage[most]{breakble-tcolorbox}

\pagestyle{plain}
\setlength{\parindent}{0pt}
\setlength{\parskip}{3pt}
\emergencystretch=2em
\hbadness=10000

\tcbset{
  proofstyle/.style={
    enhanced,
    breakable,
    boxrule=0.55pt,
    arc=1mm,
    outer arc=1mm,
    left=2.2mm,
    right=2.2mm,
    top=1.8mm,
    bottom=1.8mm,
    before skip=2.5mm,
    after skip=2.5mm,
    fonttitle=\bfseries,
  },
  outerproof/.style={
    proofstyle,
    colback=blue!3!white,
    colframe=blue!70!black,
    colbacktitle=blue!70!black,
    coltitle=white,
  },
  innerproof/.style={
    proofstyle,
    colback=green!3!white,
    colframe=green!55!black,
    colbacktitle=green!55!black,
    coltitle=white,
  },
  deepproof/.style={
    proofstyle,
    colback=orange!4!white,
    colframe=orange!80!black,
    colbacktitle=orange!80!black,
    coltitle=white,
  },
}

\newtcolorbox{outerproofbox}[1][]{outerproof,#1}
\newtcolorbox{innerproofbox}[1][]{innerproof,#1}
\newtcolorbox{deepproofbox}[1][]{deepproof,#1}

\newcommand{\algebrablock}{%
Let \(V\) be a finite-dimensional vector space and let \(f:V\to W\) be a
linear map. Choose a basis \(u_1,\dots,u_k\) of \(\ker f\), extend it to a
basis \(u_1,\dots,u_k,v_1,\dots,v_r\) of \(V\), and compare the coefficients
in that basis.

Assume that a linear relation among the image vectors is given by
\[
  a_1f(v_1)+\cdots+a_rf(v_r)=0.
\]
By linearity this is equivalent to
\[
  f(a_1v_1+\cdots+a_rv_r)=0.
\]
Therefore \(a_1v_1+\cdots+a_rv_r\in\ker f\). Since the kernel is spanned by
\(u_1,\dots,u_k\), there are coefficients \(b_1,\dots,b_k\) such that
\[
  a_1v_1+\cdots+a_rv_r=b_1u_1+\cdots+b_ku_k.
\]
Moving all terms to one side gives
\[
  b_1u_1+\cdots+b_ku_k-a_1v_1-\cdots-a_rv_r=0.
\]
Because \(u_1,\dots,u_k,v_1,\dots,v_r\) is a basis, every coefficient in this
linear combination is zero. Hence
\[
  a_1=\cdots=a_r=0.
\]
This proves that \(f(v_1),\dots,f(v_r)\) are linearly independent.

The same coefficient comparison is used again after restricting the map. If
\(U\subset V\) and an element of \(U\) is written as
\[
  x=c_1u_1+\cdots+c_ku_k+d_1v_1+\cdots+d_rv_r,
\]
then \(f(x)=0\) implies \(x\in\ker f\), so the \(v_i\)-coefficients vanish.
This sentence must appear after the displayed formula with ordinary vertical
spacing, without being pulled into the formula area.
}

\begin{document}

\begin{center}
{\Large Nested display math overlap regression}\\[2mm]
{\small A4 sample for display math inside nested breakable boxes}
\end{center}

This file checks a real regression pattern: a nested breakable fragment that
contains display math near a page break. The source does not add manual
vertical space around the displayed formulas.

\section*{Two levels with titles}

\lipsum[1-4]

\begin{outerproofbox}[title={Outer proof box}]
Text before the nested proof.

\begin{innerproofbox}[title={Inner proof box},notitle after break]
\algebrablock

\lipsum[5-8]
\end{innerproofbox}

Text after the nested proof.
\end{outerproofbox}

\section*{Three levels, no continuation title on the middle level}

\begin{outerproofbox}[title={Outer proof box},notitle after break]
Text before the titleless middle box.

\begin{innerproofbox}[notitle after break]
Text before the deep proof.

\begin{deepproofbox}[title={Deep proof box},notitle after break]
\algebrablock

\lipsum[9-12]
\end{deepproofbox}

Text after the deep proof.
\end{innerproofbox}

Text after the titleless middle box.
\end{outerproofbox}

\end{document}
